\vspace{0.8cm}
\section{Multivalued linear models}
Simple linear regression models could be extended using multiple predictors too: we switch
from a mean response that is a line to something forming a hyperplane in $\mathbb{R}^{p+1}$ (with $p$ predictors): 
\begin{center}
    $\mathbb{E}[Y|X_1=x_1,X_2=x_2]=\beta_0+\beta_1x_1+\beta_2x_2$.    
\end{center}
\begin{figure}[h]
    \centering
    \includegraphics[width=0.4\textwidth]{linearmodel/mlm/mlm_example.png}
\end{figure}\hfill\\

Trying to give an intepretation to the parameters in this new scenario, we have that
\begin{itemize}
    \item $\beta_1$ is the change in mean of $Y$ for the unitary increment of $X_1$, keeping $X_2$
    fixed
    \item $\beta_2$ is the change in mean of $Y$ for the unitary increment of $X_2$, keeping $X_1$
    fixed    
\end{itemize}
So each parameter can stil be interpreted as a slope, but slopes of a plane viewed from \diff angles. This interpretation is conditioned by the other paramters, not marginal.
\newpage
\begin{figure}[h]
    \centering
    \includegraphics[width=0.4\textwidth]{linearmodel/mlm/mlm_slope_intepret.png}
\end{figure}

So in general we have predictors $x_i\in\real^p$, responses $y_i\in\real$ and noise \seps under the assumptions $\mathcal{E}[\varepsilon|X]=0, \varepsilon\sim\mathcal{N}(0,\sigma^2I_n)$ \arrow using the matrix formulation $Y=\beta^TX+\varepsilon$.
\begin{figure}[h]
    \centering
    \includegraphics[width=0.4\textwidth]{linearmodel/mlm/matrix_not_mlm.png}
\end{figure}

To estimate the parameters we minimize the \textbf{Ordinary Least Squares}:
$$
    \hat{\beta} = \arg\min_{\beta\in\real^p} RSS(\beta) = \sum_{i=1}^n (y_i-x_i^T\beta)^2 = \left\|
    \mathbf{y}-\mathbf{X}\beta\right\|_2^2
$$
To do so we take the differential w.r.t. the parameters:
$$
    \frac{\partial}{\partial\beta} \|y-X\beta\|_2^2 = -sX^T(y-X\beta)=0 \Rightarrow \textcolor{red}{\hat{\beta} = (X^TX)^{-1}X^Ty}
$$

\subsection{Model assesment}
Under the \textit{Gaussian noise} assumption we have $\hat{\beta}\sim \mathcal{N}(\beta,\sigma^2(X^T
X)^{-1})$, giving us an estimator:
\begin{itemize}
    \item \textbf{unbiased} $\mathbb{E}[\hat{\beta}|X]=\beta$
    \item with \textbf{covariance} $Var(\hat{\beta}|X)=\sigma^2(X^TX)^{-1}$
    \item with \textbf{variance} $\hat{\sigma}^2=\cfrac{RSS}{n-p}$ (although is the same formul of simple linear regression, now we have less degree of freedom)
\end{itemize}
As for the simple linear regression, we could asses the model using the same tools:
\begin{itemize}
    \item std error $\text{SE}(\hat{\beta}_j) = \sqrt{\hat{\sigma}^2 [(X^\top X)^{-1}]_{jj}}$
    \item t-test for one parameter $H_0: \beta_j = 0, \quad t = \frac{\hat{\beta}_j}{\text{SE}
    (\hat{\beta}_j)}$
    \item confidence interval $\hat{\beta}_j \pm t_{n-p, 1-\alpha/2} \text{SE}(\hat{\beta}_j)$
\end{itemize}
Remeber that all the assesment are considered holding the other parameters fixed.

\newpage
\subsubsection{Variability explained}
\begin{figure}[h!]
    \centering
    \includegraphics[width=0.5\textwidth]{linearmodel/mlm/mlm_var_explained.png}
    \caption*{
    As we could see, the concept and metrics remain the same as for simple linear 
    regressionm, therefore the extension to multiple predictor doesn't change this part.}
\end{figure}

\subsubsection{Model usefulness}
The concept is the same as for simple linear regression: the usefulness is measured in term of how
much our model links the predictors and the responses:
\begin{center}
    $H_0:\beta_0=\beta_1=\dots\beta_{p-1}=0$ vs $H_1: \exists i:\beta_i\ne 0$ \\
    $$
        F= \cfrac{ESS/(p-1)}{RSS/(n-p)} \sim F_{p-1,n-p}
    $$
\end{center}
W.r.t. simple linear regression now the F-test is no more equivalent to the t-test due 
the fact of testing many parameters jointly.\\
\vspace{0.2cm}
We could also think to test a reduced set of parameters $p_{red}$ called \textbf{nested 
model}: to compare the reduced model with the full model
\begin{center}
    $F= \cfrac{(RSS_{full}-RSS_{red})/q}{RSS_{full}/(n-p_{full})} \sim F_{q,n-p_{full}}$ with 
    $q=p_{full}-p_{red}$
\end{center}

\subsection{Multicollinearity}
The main issue with real data is that \textit{predictors are not independent}: we have 
\textbf{multicollinearity} := one predictor is highly predictable from the others.
We have \textbf{perfect collinearity} when the relationship is linear or \textbf{near 
collinearity} when the combination is almost linear (highly correlated features).

\vspace{0.2cm}
The multicollinearity problem is not something bad in general, it's a property of the data, but it means that the predictors are redundant, with all the unpleasant implications\\

\vspace{0.2cm}
Looking at the OLS estimator $\hat{\beta}=(x^TX)^{-1}X^Ty$, we see that it requires $X^TX$ to be invertible\arrow if there is multicollinearity within the matrix there are linearly dependent columns, making the matrix non-invertible.\\
If we think about it as a linear system we have that this implies $\infty$ solutions\arrow $Xv=0$ then $X(\beta+v)=X\beta$ giving us infinite fitted values.
When the data can't distinguish between $\beta$ from $\beta+v$ is said we have 
\textbf{ill-posedness}. How do we solve the issue? Drop a variable.

\newpage
If the predictors are \textit{near-collinear} we have that $X^TX$ is still invertible, but it's \textbf{ill-conditioned}: the inverse exists but it's numerically unstable, and if we look at the variance of our parameters $Var(\hat{\beta}|X)=\sigma^2 (X^TX)^{-1}$ \arrow this leads to large std error, wide 
confidence intervals ...\\
\begin{figure}[h]
    \centering
    \includegraphics[width=0.5\textwidth]{linearmodel/mlm/var_inflation.png}
\end{figure}

However, we could have fitted values $\hat{y}=X\hat{\beta}$ pretty stable even with unstable params 
\arrow multicollinearity is much worse for intepretation/inference rather than for prediction (almost ever).

\paragraph{Variance Inflation Factor (VIF)}\hfill\\
How can we understand if some predictor might cause \textit{variance inflations}?

\vspace{0.2cm}
Take predictor $X_j$ and build a model that predicts $X_j$ from all the other predictor and consider this $R_j$, we have that
\begin{align*}
& \text{Var}(\hat{\beta}_j | X) = \frac{\sigma^2}{\sum_{i=1}^{n} (x_{ij} - \bar{x}_j)^2} \cdot \frac{1}{1 - R_j^2}\\
& \Rightarrow \textbf{VIF} = \frac{1}{1-R^2_j}
\end{align*}
If $R_j^2\rightarrow0\Rightarrow \text{VIF}\rightarrow\infty$ we see that the variance of the coefficient explodes, impling correlation between some predictors.

\subsubsection{How to deal with multicollinearity}
\begin{itemize}
    \item remove redundant predictors
    \item merge correlated predictors into one, changing its interpretability
    \item collect more data, hoping that by adding more variation we might be able to reduce colllinearity
    \item regularize (see later)
    \item apply PCA \arrow PCA regression (making the predictors independent by default)
\end{itemize}

\newpage
\subsection{Checking model assumptions}
Up to now we've made some assumptions about the error: $\mathcal{E}[\varepsilon|X]=0, Var(\varepsilon|
X)=\sigma^2, Cov(\varepsilon_i,\varepsilon_j|X)=0$ and on other aspects  of the model, but what if
this assumption doesn't hold?

\paragraph{Model linearity}
\begin{center}
    \textit{The model is linear in the coefficients} $\mathbb{E}[Y_i|X]=\beta_0+\beta_1x_{i1}
    +\beta_2x_{i2}$
\end{center}
\begin{figure}[h]
    \centering
    \includegraphics[width=0.5\textwidth]{linearmodel/mlm/lin_assumption.png}
\end{figure}
This doesn't mean that the predictors have to be linear, only the combination of predictors and parameters must be.

\paragraph{Null error mean}
\begin{center}
    $\expect[\varepsilon_i|X]=0$
\end{center}
meaning that there is no leftover information in the error and the errors are uncorrelated with the included predictors.\\
If this this assumption doesn't hold, we might have: variable bias, measurement errors in
predictors, bad model structure.

% if you figure out how to do it, do it
\begin{comment}
\begin{wrapfigure}[4]{l}{0.4\linewidth}
    \includegraphics[width=0.9\linewidth]{linearmodel/mlm/homo_residual.png}
\end{wrapfigure}
\end{comment}

\paragraph{Homoscedasticity of residuals}
\begin{center}
    \textit{The error variance should be constant} $Var(\varepsilon_i|X)=\sigma^2$
\end{center} 
This means that the error should be spread roughly the same across the fitted model. If not, the 
OLS are still unbiased but the std may be wrong\arrow t-tests and CI might be misleading.
\begin{figure}[h!]
    \centering
    \includegraphics[width=0.3\linewidth]{linearmodel/mlm/homo_residual.png}
\end{figure}

\paragraph{Independence of errors}
\begin{center}
    \textit{Errors are independent} $Cov(\varepsilon_i,\varepsilon_j|X)=0\ \forall i\ne j$
\end{center}
If errors are correlated we get distorted std errors, inference might be wrong and the actual 
information within the datset is less than the sample size.\\
This kind of assumption often fails in a lot of realistic scenarios (time series, repeated
measures ...).

\paragraph{Independence of errors}
\begin{center}
    $\varepsilon_i \sim \mathcal{N(0,\sigma^2)}$
\end{center}
This assumption is not required for the OLS estimator rather than the simplicity of having
exact t-test, exact F-test, exact CI and exact PI. In larger samples mild deviations from
normal distribution are not problematic due CLT asymptotic result. 

\subsubsection{Leverage, outliers and influential points}

\paragraph{Leverage points}\hfill\\
Leverage points are those point extremely unusual in the predictors dimension. If we take a closer 
look to our prediction we have that $\hat{y}=X\hat{\beta} \xrightarrow[OLS]{} \hat{y}=X(X^TX)^{-1}
X^Ty=Hy$ where $H$ := \textbf{hat matix} \arrow $h_{ii}$ gives the \textbf{leverage} of the 
observations, in other words how unusual the observation is in the predictor space.\\ 
\vspace{0.2cm}
high leverage \arrow unusual combination\\
low leverage \arrow typical observation\\
\begin{figure}[h!]
    \centering
    \includegraphics[width=0.6\textwidth]{linearmodel/mlm/leverage_point.png}
\end{figure}
High-leverage points following the regression model doesn't distort the model (they follow the 
exploited structure) but due them the regression line is strongly supported for extreme $x$, 
reducing the uncertanty in that region even if the information about it isn't that much, making CIs 
artificially tighter\arrow due them the model gives a false sense of generalization.

\newpage
\paragraph{Outliers}\hfill\\
Outliers are those observations that are unusual in the response, characterized by large residual
values $e_i=y_i-\hat{y}_i$. This can cause the increase of RSS, inflate the estimated noise level,
distort std error (with consequences on test and CI). Another thing is that this type of point 
might suggest possible data issues or model miss-specification.
\begin{figure}[h!]
    \centering
    \includegraphics[width=0.6\textwidth]{linearmodel/mlm/outlier.png}
\end{figure}

\paragraph{Influential points}\hfill\\
\begin{center}
    high leverage + large residual \arrow influential point
\end{center}
This type of points might have big influence on the model, causing changes in the estimation of the
coefficients, distorted predictions and inference.
\begin{figure}[h!]
    \centering
    \includegraphics[width=0.6\textwidth]{linearmodel/mlm/influential_point.png}
\end{figure}

\newpage
\subsubsection{Residual vs fitted diagnostic plots}
\begin{figure}[h!]
    \centering
    \includegraphics[width=0.8\textwidth]{linearmodel/mlm/resisual_vs_fitted_plot.png}
\end{figure}
Useful to check linearity of the mean, homoscedacity and possible outliers/leverage points.\\
\vspace{0.2cm}
\textcolor{green}{Good pattern}: andom cloud centered around zero, without visible structure.\\
\textcolor{red}{Bad pattern}: curved pattern (suggests nonlinearity), funnel shape (suggests 
heteroscedasticity), isolated points (possible outliers)\\

\subsubsection{Q-Q diagnostic plots}
\begin{figure}[h!]
    \centering
    \includegraphics[width=0.8\textwidth]{linearmodel/mlm/qq_plot.png}
\end{figure}
Usefull to check whether residuals are approximately normal or not.\\
\vspace{0.2cm}
\textcolor{green}{Good pattern}: points close to the reference line, especially in the center.\\
\textcolor{red}{Bad pattern}: strong tail departures (extreme residuals), systematic asymmetry 
(skewness), s-shaped pattern (mismatch).\\

\subsubsection{Shapiro-Wlik test}
Standard test to check whether resisuals are normally distributed or not, to be used together with
Q-Q plot to strengthen it.
\begin{center}
    \textit{$H_0$: residuals are normal $\ H_1$: residuals are not normal}
\end{center}
large p-value \arrow no evidence against normality\\
small p-value \arrow residuals probably are not normally distributed.

\newpage
\subsubsection{Scale-location plots}
\begin{figure}[h!]
    \centering
    \includegraphics[width=0.8\textwidth]{linearmodel/mlm/scale-location_plot.png}
\end{figure}
Used to check for homoscedacity.\\
\vspace{0.2cm}
\textcolor{green}{Good pattern}: roughly horizontal band with uniform spread.\\
\textcolor{red}{Bad pattern}: increasing trend (variance grows with fitted values) or decreasing 
trend.

\vspace{0.4cm}
\subsubsection{Leverage vs studentized residuals}
Used to separate unusual $X$-values from unusual $Y$/values (in other words useful to highlight
leverage, outliers and influential points.)
\begin{figure}[h!]
    \centering
    \includegraphics[width=0.7\textwidth]{linearmodel/mlm/lev_vs_stud_plot.png}
\end{figure}
